% !TEX program = lualatex
\documentclass[professionalfonts]{beamer}
\newcommand{\slidemode}{1}  % カラー投影モード
\usepackage{beamer_template}
\usepackage{enumerate}

\newcommand{\LectureNum}{12}
\newcommand{\LectureTitle}{フーリエ変換と積分定理}
\newcommand{\CourseName}{応用数学}
\renewcommand{\TermName}{前期}

\begin{document}

\begin{frame}{第12回　フーリエ変換と積分定理}
  \begin{block}{今回の内容}
  フーリエ変換と積分定理
  \end{block}
  \vfill\hfill{\scriptsize \textcolor{gray}{第3章 フーリエ解析　§2}}
\end{frame}

\begin{frame}{定義：フーリエ変換}
  \begin{block}{}
  \medskip\textbf{フーリエ変換}　$\bigl(F(u) = F[f(x)]\bigr)$：
  \[
    F(u) = \int_{-\infty}^{\infty} f(x) e^{-iux} dx
  \]
  \medskip\textbf{逆フーリエ変換}　$\bigl(f(x) = F^{-1}[F(u)]\bigr)$：
  \[
    f(x) = (1/2\pi) \int_{-\infty}^{\infty} F(u) e^{iux} du
  \]
  {\footnotesize \textcolor{gray}{$F(u)$ を $f(x)$ のフーリエ変換、 $f(x)$ を $F(u)$ の逆フーリエ変換という}}
  \end{block}
\end{frame}

\begin{frame}{定理：フーリエ積分定理}
  \begin{block}{}
  \textbf{条件}：
  \begin{itemize}
    \item $f(x)$ は任意の有限区間で区分的に滑らか
    \item $\int_{-\infty}^{\infty} |f(x)| dx < \infty$（絶対可積分）
  \end{itemize}
  \[
    f(x) = (1/\pi) \int_{0}^\infty [\int_{-\infty}^{\infty} f(\tau)\cos(u(x-\tau)) d\tau] du
  \]
  \end{block}
\end{frame}

\begin{frame}{フーリエ余弦変換}
  \begin{block}{}
  \[
    Fc(u) = \int_{0}^\infty f(x)\cos(ux) dx
  \]
  \[
    f(x) = (2/\pi) \int_{0}^\infty Fc(u)\cos(ux) du
  \]
  {\footnotesize $f(x)$ が偶関数の場合に対応}
  \end{block}
\end{frame}

\begin{frame}{フーリエサイン変換}
  \begin{block}{}
  \[
    Fs(u) = \int_{0}^\infty f(x)\sin(ux) dx
  \]
  \[
    f(x) = (2/\pi) \int_{0}^\infty Fs(u)\sin(ux) du
  \]
  {\footnotesize $f(x)$ が奇関数の場合に対応}
  \end{block}
\end{frame}

\begin{frame}{例題1　（p.\,91）}
  \begin{exampleblock}{問題}
    $f(x) = e^{-|x|}$ のフーリエ変換を求めよ\\[2pt]
  \end{exampleblock}
\end{frame}

\begin{frame}{例題1　解答}
  $F(u) = \int_{-\infty}^{\infty} e^{-|x|} e^{-iux} dx$\\[4pt]
  $= \int_{-\infty}^{0} e^{x} e^{-iux} dx + \int_{0}^{\infty} e^{-x} e^{-iux} dx$\\[4pt]
  $= 1/(1-iu) + 1/(1+iu)$
  \\[8pt]\textcolor{red}{\textbf{答}}：$F(u) = 2/(1+u^{2})$
  {\footnotesize \textcolor{gray}{\textit{$x\to-\infty$ で $e^{(1-iu)x}\to 0$、$x\to+\infty$ で $e^{-(1+iu)x}\to 0$（実部 $e^x, e^{-x}$ より）}}}
\end{frame}

\begin{frame}{演習問題}
  \begin{exampleblock}{自分で解いてみよう}
  \begin{description}
    \item[問1] 次の関数のフーリエ変換を求めよ
    $(1) f(x) = \begin{cases} 0 & (x > 0) \\ e^x & (x \leq 0) \end{cases}$\\
    $(2) g(x) = \begin{cases} xe^{-x} & (x > 0) \\ 0 & (x \leq 0) \end{cases}$\\
  \end{description}
  \end{exampleblock}
\end{frame}

\begin{frame}{問1　解答}
  $(1)\; F(u) = \displaystyle\int_{-\infty}^{0} e^x e^{-iux}\,dx = \int_{-\infty}^{0} e^{(1-iu)x}\,dx$\\[4pt]
  $= \left[\dfrac{e^{(1-iu)x}}{1-iu}\right]_{-\infty}^{0} = \dfrac{1}{1-iu} \quad (\Re(1-iu)=1>0)$\\[6pt]
  $(2)\; G(u) = \displaystyle\int_{0}^{\infty} x e^{-x} e^{-iux}\,dx = \int_{0}^{\infty} x e^{-(1+iu)x}\,dx$\\[4pt]
  $\displaystyle\int_{0}^\infty x e^{-ax}\,dx = \dfrac{1}{a^{2}}$ より $G(u) = \dfrac{1}{(1+iu)^{2}}$
  \\[8pt]\textcolor{red}{\textbf{答}}：
  $(1)\; F(u) = \dfrac{1}{1-iu}$\qquad
  $(2)\; G(u) = \dfrac{1}{(1+iu)^{2}}$\\
\end{frame}

\begin{frame}{例題2　（p.\,92）}
  \begin{exampleblock}{問題}
    $f(x) = \begin{cases} 1 & (|x| \leq 1) \\ 0 & (|x| > 1) \end{cases}$ のフーリエ変換を求めよ\\[2pt]
  \end{exampleblock}
\end{frame}

\begin{frame}{例題2　解答}
  $F(u) = \int_{-1}^{1} e^{-iux} dx = [-e^{-iux}/(iu)]_{-1}^{1}$\\[4pt]
  $= (e^{iu} - e^{-iu})/(iu) = 2\sin u / u$
  \\[8pt]\textcolor{red}{\textbf{答}}：$F(u) = 2\sin u / u$
\end{frame}

\begin{frame}{演習問題}
  \begin{exampleblock}{自分で解いてみよう}
  \begin{description}
    \item[問2] 次の関数のフーリエ変換を求めよ
    $(1) f(x) = \begin{cases} 1 & (0 \leq x \leq 1) \\ 0 & (x < 0, x > 1) \end{cases}$\\
    $(2) g(x) = \begin{cases} x & (|x| \leq 1) \\ 0 & (|x| > 1) \end{cases}$\\
    \item[問3] 問2(1)の関数に積分定理を適用して、次の等式を証明せよ
    $(1) (1/2\pi)\int_{-\infty}^{\infty} (1-e^{-iu})/iu \cdot e^{iux} du = \begin{cases} 1 & (0<x<1) \\ 1/2 & (x=0,1) \\ 0 & (x<0, x>1) \end{cases}$\\
    $(2) \int_{-\infty}^{\infty} \sin u / u  du = \pi$\\
  \end{description}
  \end{exampleblock}
\end{frame}

\begin{frame}{問2(1)　解答}
  $(1)\; F(u) = \displaystyle\int_{0}^1 e^{-iux}\,dx = \left[\dfrac{e^{-iux}}{-iu}\right]_{0}^1 = \dfrac{1-e^{-iu}}{iu}$
  \\[8pt]\textcolor{red}{\textbf{答}}：$(1)\; F(u) = \dfrac{1-e^{-iu}}{iu}$\\
\end{frame}

\begin{frame}{問2(2)　解答}
  $(2)\; g(x)$ は奇関数 $:\; G(u) = \displaystyle\int_{-1}^{1} x e^{-iux}\,dx = -2i\int_{0}^1 x \sin(ux)\,dx$\\[4pt]
  $\displaystyle\int_{0}^1 x \sin(ux)\,dx = \left[-\dfrac{x \cos(ux)}{u}\right]_{0}^1 + \dfrac{1}{u}\int_{0}^1 \cos(ux)\,dx = -\dfrac{\cos u}{u} + \dfrac{\sin u}{u^{2}}$\\[6pt]
  $G(u) = -2i\!\left(-\dfrac{\cos u}{u} + \dfrac{\sin u}{u^{2}}\right) = \dfrac{2i\cos u}{u} - \dfrac{2i\sin u}{u^{2}}$
  \\[8pt]\textcolor{red}{\textbf{答}}：$(2)\; G(u) = \dfrac{2i\cos u}{u} - \dfrac{2i\sin u}{u^{2}}$\\
\end{frame}

\begin{frame}{問3(1)　解答}
  $(1)$ フーリエ積分定理 $:\; f(x) = \dfrac{1}{2\pi}\displaystyle\int_{-\infty}^{\infty} F(u)\,e^{iux}\,du$\\[6pt]
  連続点では $f(x)$、不連続点 $x=0,1$ では $\dfrac{f(x+0)+f(x-0)}{2} = \dfrac{1}{2}$
  \\[8pt]\textcolor{red}{\textbf{答}}：証明終
\end{frame}

\begin{frame}{問3(2)　解答}
  $(2)\; x=1$（不連続点）を代入 $:\; \dfrac{1}{2} = \dfrac{1}{2\pi}\displaystyle\int_{-\infty}^{\infty}\dfrac{1-e^{-iu}}{iu}\cdot e^{iu}\,du$\\[6pt]
  $\dfrac{(1-e^{-iu})e^{iu}}{iu} = \dfrac{e^{iu}-1}{iu} = \dfrac{\sin u}{u} + \dfrac{i(1-\cos u)}{u}$\\[6pt]
  実部（偶関数）のみ寄与 $:\; \dfrac{1}{2} = \dfrac{1}{2\pi}\displaystyle\int_{-\infty}^{\infty} \dfrac{\sin u}{u}\,du$\\[6pt]
  $\therefore\; \displaystyle\int_{-\infty}^{\infty} \dfrac{\sin u}{u}\,du = \pi$
  \\[8pt]\textcolor{red}{\textbf{答}}：証明終
\end{frame}

\begin{frame}{例題3　（p.\,95）}
  \begin{exampleblock}{問題}
    $f(x) = \begin{cases} 1-|x| & (|x| \leq 1) \\ 0 & (|x| > 1) \end{cases}$ のフーリエ余弦変換を求めよ\\[2pt]
  \end{exampleblock}
\end{frame}

\begin{frame}{例題3　解答}
  $f(x)$ は偶関数なのでフーリエ余弦変換を用いる\\[4pt]
  $Fc(u) = \int_{0}^{1} (1-x) \cos(ux) dx$\\[4pt]
  $= [(1-x)\sin(ux)/u]_{0}^1 + (1/u)\int_{0}^1 \sin(ux) dx$\\[4pt]
  $= 0 + (1/u)[-\cos(ux)/u]_{0}^1 = (1-\cos u)/u^{2}$
  \\[8pt]\textcolor{red}{\textbf{答}}：$Fc(u) = (1-\cos u)/u^{2}$
\end{frame}

\begin{frame}{演習問題}
  \begin{exampleblock}{自分で解いてみよう}
  \begin{description}
    \item[問4] 次の関数のフーリエサイン変換を求めよ
    $f(x) = \begin{cases} -1 & (-1 \leq x < 0) \\ 1 & (0 \leq x \leq 1) \\ 0 & (|x| > 1) \end{cases}$\\
  \end{description}
  \end{exampleblock}
\end{frame}

\begin{frame}{問4　解答}
  $f(x)$ は奇関数なのでフーリエサイン変換を用いる\\[6pt]
  $Fs(u) = \displaystyle\int_{0}^1 \sin(ux)\,dx = \left[-\dfrac{\cos(ux)}{u}\right]_{0}^1 = \dfrac{1-\cos u}{u}$
  \\[8pt]\textcolor{red}{\textbf{答}}：$Fs(u) = \dfrac{1-\cos u}{u}$
\end{frame}

\end{document}
